%=======================================================================
%  Appendix_proof_sketches.tex
%
%  Companion note to the ContEst paper.  It collects *sketch* proofs of the
%  two regularity propositions that appear in the paper's appendices:
%
%     Appendix A  ==  Proposition (LQG regularity)      [prop:lqgreg]
%     Appendix B  ==  Proposition (fixed-gamma H_infty regularity) [prop:hinfreg]
%
%  Each sketch states the claim, then gives the argument at the level of the
%  key identities and the theorems they rest on -- naming those theorems and
%  explaining in one or two lines *why* they apply -- rather than the fully
%  expanded computation.  The complete LQG argument (with every Schur
%  complement written out) is in Thm1_proof.tex in the same folder; this note
%  is the skimmable version and additionally covers the H_infty case.
%
%  The document compiles on its own (\documentclass{article}) and restates
%  every object it needs, so no cross-file references are required.
%=======================================================================
\documentclass[11pt]{article}

\usepackage[margin=1in]{geometry}
\usepackage{amsmath,amssymb,amsthm}
\usepackage{mathtools}
\usepackage{bm}
\usepackage{hyperref}

\newtheorem{proposition}{Proposition}
\newtheorem{lemma}{Lemma}
\theoremstyle{definition}
\newtheorem{remark}{Remark}

\DeclareMathOperator{\tr}{tr}
\DeclareMathOperator{\rank}{rank}
\DeclareMathOperator{\range}{range}
\DeclareMathOperator{\diag}{diag}
\newcommand{\R}{\mathbb{R}}
\newcommand{\Sym}{\mathbb{S}}
\newcommand{\Acl}{A_{\mathrm{cl}}}
\newcommand{\Th}{\theta}
\newcommand{\Nbhd}{\Theta_{0}}
\newcommand{\half}{{1/2}}
\newcommand{\inner}[2]{\langle #1,\, #2\rangle}
\newcommand{\Pe}{P^{\varepsilon}}
\newcommand{\Vwc}{V_{\mathrm{wc}}}
\newcommand{\Vest}{V_{\mathrm{est}}}

\title{\bf Sketch proofs of the ContEst appendix regularity propositions}
\author{ContEst supplementary note}
\date{\today}

\begin{document}
\maketitle

\begin{abstract}
This note gives sketch proofs of the two regularity results proved in the
appendices of the ContEst paper: the LQG regularity proposition
(Appendix~A), under which the covariance-SDP design gradients of the LQG inner
problem are exact, and the fixed-$\gamma$ $H_\infty$ regularity proposition
(Appendix~B), its robust twin for the game-Riccati inner problem. For each we
state the claim, then argue at the level of the decisive identities and the
classical theorems they invoke --- Slater/conic duality, the discrete algebraic
Riccati equation (DARE) and its game analog (GARE), the implicit function theorem
via the Stein operator, the Danskin/Bonnans--Shapiro value-derivative formula, and
the Alizadeh--Haeberly--Overton characterization of nondegeneracy --- explaining
briefly why each applies. The fully expanded LQG argument is in the companion
\texttt{Thm1\_proof.tex}; this note is the compact version and additionally covers
$H_\infty$.
\end{abstract}

%=======================================================================
\section{Common setup and notation}\label{sec:setup}
%=======================================================================

Fix a design point $\Th$ in an open set $\Nbhd\subseteq\R^{r_\theta}$ and the
discrete-time model
\begin{equation}\label{eq:model}
x_{k+1}=A_\Th x_k+B_\Th u_k+w_k,\qquad y_k=C_\Th x_k+v_k,
\end{equation}
$x_k\in\R^{n}$, $u_k\in\R^{m}$, $y_k\in\R^{p}$, with $C^1$ data
$\Th\mapsto(A_\Th,B_\Th,C_\Th,W,V,Q,R)$, noise covariances $W,V$ and stage weights
$Q,R$. We write $\Sym^q$ for symmetric $q\times q$ matrices,
$\inner{X}{Y}=\tr(X^\top Y)$, and call $M$ \emph{Schur} if $\rho(M)<1$. Both inner
problems are instances of the paper's generic convex program
\begin{equation}\label{eq:generic}
V(\Th)=\min_{z}\ \phi(z,\Th)\ \ \text{s.t.}\ \ g(z,\Th)=0,\ G(z,\Th)\in\mathcal K,
\end{equation}
with $\mathcal K$ the PSD cone, Lagrangian $L=\phi+\inner{\mu}{g}+\inner{S}{G}$,
optimal triple $(z^\star,\mu^\star,S^\star)$. The paper's Theorem~1 (envelope
theorem) states that under its Assumption~1,
\begin{equation}\label{eq:envelope}
\nabla_\Th V=\partial_\Th\phi+\inner{\mu^\star}{\partial_\Th g}
+\inner{S^\star}{\partial_\Th G},
\end{equation}
the partials taken with $(z^\star,\mu^\star,S^\star)$ frozen. Assumption~1 asks,
on $\Nbhd$: (A1) convexity, $C^1$ data, Slater; (A2) unique primal minimizer;
(A3) nondegeneracy/strict complementarity, hence a unique dual; (A4) $\Th$ enters
only through finitely many $C^1$ data objects. The two propositions below reduce
(A1)--(A4) to \emph{system-theoretic} hypotheses that are checkable from the model.

Throughout, the standing system-theoretic hypotheses are
\begin{itemize}
\item[\textup{(S1)}] the data maps are $C^1$ in $\Th$;
\item[\textup{(S2)}] $(A_\Th,B_\Th)$ stabilizable and $(A_\Th,C_\Th)$ detectable;
\item[\textup{(S3)}] $W,V,R\succ0$, $Q\succeq0$ with $(A_\Th,Q^{\half})$ detectable
      (automatic when $Q\succ0$).
\end{itemize}

%=======================================================================
\section{Appendix A: LQG regularity --- sketch}\label{sec:lqg}
%=======================================================================

The LQG inner value splits, by the fixed-$\Th$ separation principle, into a
control SDP and an estimation SDP. In covariance form the control SDP is
\begin{equation}\label{eq:Jcont}
J_{\mathrm{cont}}(\Th)=\min_{\Sigma,L,Z_0}\ \tr(Q\Sigma)+\tr(RZ_0)\ \
\text{s.t.}\ \ \Sigma\succ0,\
\begin{bsmallmatrix}Z_0&L\\L^\top&\Sigma\end{bsmallmatrix}\succeq0,\
\mathcal G_{\mathrm c}\succeq0,
\end{equation}
with $\mathcal G_{\mathrm c}=
\begin{bsmallmatrix}\Sigma-W & A_\Th\Sigma+B_\Th L\\
(A_\Th\Sigma+B_\Th L)^\top & \Sigma\end{bsmallmatrix}$ and optimal gain
$K^\star=L^\star(\Sigma^\star)^{-1}$; the estimation SDP is its transpose dual
(Section~\ref{sec:lqg-est}).

\begin{proposition}[LQG regularity]\label{prop:lqg}
Under \textup{(S1)}--\textup{(S3)}, for \eqref{eq:Jcont} and its estimation dual and
every $\Th\in\Nbhd$: \textup{(a)} Slater holds; \textup{(b)} the primal minimizer is
unique and equals the stabilizing DARE solution, $\Sigma^\star\succ0$;
\textup{(c)} the value is $C^1$ and, for \emph{every} optimal dual, \eqref{eq:envelope}
returns $\nabla_\Th V$ (so the design gradients read from the solver's duals are
exact); \textup{(d)} if in addition the solution is strictly complementary and
primal-nondegenerate, the dual is unique and Assumption~1 holds verbatim. Moreover
the $(1,1)$ dual block of $\mathcal G_{\mathrm c}$ equals the control Riccati
solution, $S_{11}^\star=P$.
\end{proposition}

\paragraph{(a) Slater from stabilizability/detectability.}
Pick $K$ with $\Acl:=A_\Th+B_\Th K$ Schur (exists by (S2)). For any $\varepsilon>0$
the discrete Lyapunov equation $\Sigma=\Acl\Sigma\Acl^\top+W+\varepsilon I$ has a
unique solution $\Sigma=\sum_{k\ge0}\Acl^k(W+\varepsilon I)\Acl^{\top k}\succ0$
because $\Acl$ is Schur and $W+\varepsilon I\succ0$ (\emph{Lyapunov stability
theorem} for discrete systems). Setting $L=K\Sigma$, $Z_0=K\Sigma K^\top+\varepsilon I$,
the Schur complements of the three LMIs are $\varepsilon I\succ0$, so the point is
\emph{strictly} feasible. Slater for a bounded-below convex conic program gives
zero duality gap and dual attainment by the \emph{conic strong-duality theorem}
\cite[Thm.~2.4.1]{BonnansShapiro},\cite[\S5.9]{BoydVandenberghe}. This is (A1); (A4)
is immediate since $\Th$ enters only through $(A_\Th,B_\Th,W,Q,R)$, all $C^1$.

\paragraph{(b) Unique primal $=$ DARE solution.}
Two classical facts do the work. First, the covariance SDP is an exact convex lift
of the stationary LQR problem: for feasible $(\Sigma,L,Z_0)$ with $K=L\Sigma^{-1}$,
the Schur complement of $\mathcal G_{\mathrm c}$ reads
$\Sigma\succeq\Acl\Sigma\Acl^\top+W$, and $W\succ0$ forces $\Acl$ Schur (a
\emph{Lyapunov-inequality / eigenvector} argument \cite[Ch.~5]{BoydLMI}); iterating
the resulting Stein inequality gives $\Sigma\succeq\Sigma_K$, so the objective
dominates the stationary LQR cost $J(K)=\tr(Q\Sigma_K)+\tr(RK\Sigma_KK^\top)$.
Second, the \emph{completion-of-squares identity}
\begin{equation}\label{eq:cos}
J(K)-\tr(PW)=\tr\!\big((K-K^\star)^\top(R+B_\Th^\top PB_\Th)(K-K^\star)\Sigma_K\big),
\end{equation}
with $P$ the stabilizing solution of the DARE
$P=A_\Th^\top PA_\Th-A_\Th^\top PB_\Th(R+B_\Th^\top PB_\Th)^{-1}B_\Th^\top PA_\Th+Q$,
is $\ge0$ and vanishes iff $K=K^\star$ (here $R+B_\Th^\top PB_\Th\succ0$,
$\Sigma_K\succ0$). Existence/uniqueness of the stabilizing $P\succeq0$ is the
\emph{DARE theorem} under (S2)--(S3) \cite[Ch.~14--15]{LancasterRodman},%
\cite[Ch.~6]{AndersonMoore}. Hence the minimizer $\Sigma^\star=\Sigma_{K^\star}
\succ0$, $L^\star=K^\star\Sigma^\star$ is unique and $J_{\mathrm{cont}}=\tr(PW)$,
which is (A2). (For $Q\succeq0$ the final pinning of $\Sigma^\star$ uses
$(A_\Th,Q^{\half})$ detectability; for $Q\succ0$ it is immediate.)

\paragraph{(c) $C^1$ value and dual-invariant gradient (the operational core).}
Because $J_{\mathrm{cont}}=\tr(PW)$ and $P(\Th)$ is a $C^1$ function of the data,
the value is $C^1$: apply the \emph{implicit function theorem} to the DARE residual
$\mathcal R(P,\Th)=0$, whose $P$-derivative at a stabilizing root is the Stein
operator $\delta P\mapsto\delta P-\Acl^{\star\top}\delta P\,\Acl^\star$, invertible
exactly because $\Acl^\star$ is Schur (eigenvalues $1-\lambda_i\bar\lambda_j\neq0$).
Now the key point that makes the gradient safe to read off \emph{any} returned dual:
by the \emph{Danskin/Bonnans--Shapiro directional-derivative theorem} for optimal
values of convex programs with a unique primal solution
\cite[Thm.~4.24]{BonnansShapiro},
\[
V'(\Th;d)=\min_{(\mu,S)\in\mathcal D^\star(\Th)}
\inner{\nabla_\Th L(z^\star,\mu,S,\Th)}{d},
\]
where $\mathcal D^\star$ is the (compact convex) dual optimal set. Since $V$ is
$C^1$, the left side is \emph{linear} in $d$, so the min equals
$\inner{\nabla_\Th V}{d}$ for all $d$; a fixed dual gives
$\inner{\nabla_\Th L}{d}\ge\inner{\nabla_\Th V}{d}$ for all $d$ (including $-d$),
forcing equality $\nabla_\Th L=\nabla_\Th V$. Thus \emph{every} optimal dual yields
the same, correct gradient, and Corollary~1 of the paper holds under
(S1)--(S3) \emph{without} strict complementarity. Specialized to the Lyapunov LMI,
$\nabla_\Th L$ is the paper's $\partial J_{\mathrm{cont}}/\partial A_\Th=-2S_{12}
\Sigma^\star$, etc.

\paragraph{The dual identity.}
Complementary slackness aligns $\range S^\star$ with
$\ker\mathcal G_{\mathrm c}^\star=\range[\,I;\,-\Acl^{\star\top}\,]$ (the tight
Lyapunov block has rank $n$), so $S^\star=[\,I;-\Acl^{\star\top}\,]M[\,I;-\Acl^{\star
\top}\,]^\top$ and $S_{12}^\star=-S_{11}^\star\Acl^\star$. Lagrangian stationarity in
$\Sigma$ then gives the closed-loop \emph{Stein equation}
$S_{11}=Q+K^{\star\top}RK^\star+\Acl^{\star\top}S_{11}\Acl^\star$, whose unique
solution ($\Acl^\star$ Schur) is $P$; hence $S_{11}^\star=P$ and, via
$S_{12}^\star=-P\Acl^\star$, $\partial J_{\mathrm{cont}}/\partial A_\Th=-2S_{12}^\star
\Sigma^\star=2P\Acl^\star\Sigma^\star$, matching the Riccati gradient. (This is the
identity corrected in the paper's Appendix A: the off-diagonal dual block carries
$\Acl^\star$, not $K^{\star\top}$.)

\paragraph{(d) Strict complementarity $\Rightarrow$ unique dual.}
Primal nondegeneracy at $z^\star$ implies the dual optimizer is unique and strict
complementarity implies the primal--dual solution map is differentiable, by the
\emph{Alizadeh--Haeberly--Overton / Shapiro} characterization for SDPs
\cite[Thms.~10,15]{AHO},\cite[Thm.~4.1]{Shapiro97}. Here $\rank\mathcal
G_{\mathrm c}^\star=n$ always, so strict complementarity reduces to $\rank
S^\star=n$ with $S^\star\mathcal G_{\mathrm c}^\star=0$, guaranteed by rank counting
from $P\succeq Q\succ0$, $R\succ0$. This supplies (A3), completing Assumption~1.

\subsection{The estimation SDP by transpose duality}\label{sec:lqg-est}
The estimation SDP is \eqref{eq:Jcont} under the relabeling
$A_\Th\mapsto A_\Th^\top$, $B_\Th\mapsto C_\Th^\top$, $W\mapsto Q$, $Q\mapsto W$,
$R\mapsto V$, $(\Sigma,L,Z_0)\mapsto(\Pe,-F,Z_1)$. The hypotheses map correctly:
$(A_\Th,B_\Th)$ stabilizable $\leftrightarrow(A_\Th,C_\Th)$ detectable, and
$(A_\Th,Q^{\half})$ detectable $\leftrightarrow(A_\Th,W^{\half})$ stabilizable
(automatic since $W\succ0$ makes it controllable). Every step above is invariant
under this relabeling, so all conclusions transfer, giving the unique
$\Pe{}^\star=\Sigma^\varepsilon$ (Kalman error covariance) and $S_{11}^{\prime\star}=
\Sigma^\varepsilon$.

%=======================================================================
\section{Appendix B: fixed-$\gamma$ $H_\infty$ regularity --- sketch}\label{sec:hinf}
%=======================================================================

For the robust inner problem the paper uses, at a \emph{fixed} attenuation
$\gamma$, the game Riccati route. With disturbance channel $E_\Th=W_\Th^{\half}$,
augmented input $\tilde B_\Th=[\,B_\Th\ E_\Th\,]$ and \emph{indefinite} game weight
$\tilde R=\diag(R,-\gamma^2I)$, the guaranteed cost is
$\Vwc(\Th)=\tr(XW_\Th)$ with $X\succeq0$ the stabilizing solution of the game GARE
\begin{equation}\label{eq:gare}
X=Q+A_\Th^\top XA_\Th-A_\Th^\top X\tilde B_\Th
(\tilde R+\tilde B_\Th^\top X\tilde B_\Th)^{-1}\tilde B_\Th^\top XA_\Th,
\end{equation}
saddle gain $\tilde K=-(\tilde R+\tilde B_\Th^\top X\tilde B_\Th)^{-1}
\tilde B_\Th^\top XA_\Th$, closed loop $\Acl=A_\Th+\tilde B_\Th\tilde K$.

\begin{proposition}[Fixed-$\gamma$ $H_\infty$ regularity]\label{prop:hinf}
Fix $\gamma$ above the optimal attenuation and let $X\succeq0$ solve \eqref{eq:gare}
with $R+B_\Th^\top XB_\Th\succ0$ and $\Acl$ Schur. Then $X$ is a locally unique
$C^1$ function of $\Th$, $\Vwc(\Th)=\tr(XW_\Th)\in C^1$, with the closed-form
gradient \eqref{eq:garegrad} below. The dual filter GARE is regular under the
analogous conditions, giving $\Vest(\Th)=\tr(M\Sigma_e)\in C^1$.
\end{proposition}

\noindent The proof is the Riccati (implicit-function) half of
Section~\ref{sec:lqg}, carried through with an indefinite weight; there is no SDP
dual and no strict-complementarity requirement.

\paragraph{(a) Stabilizing solution.}
For $\gamma$ above the optimal attenuation $\gamma^\star$, the \emph{discrete
dynamic-game (bounded-real) Riccati theorem} gives a unique $X\succeq0$ solving
\eqref{eq:gare} with $R+B_\Th^\top XB_\Th\succ0$, full pivot
$\tilde R+\tilde B_\Th^\top X\tilde B_\Th$ nonsingular of inertia $(m,r_w)$, and
$\Acl$ Schur \cite{BasarBernhard},\cite[Ch.~2]{BoydLMI}. This is the fixed-$\gamma$
analog of the DARE existence used in Section~\ref{sec:lqg}; approaching
$\gamma^\star$ makes the pivot lose rank (the minimal-$\gamma$ SDP degeneracy the
paper avoids), so we stay strictly above it.

\paragraph{(b) Lyapunov (completion-of-squares) form.}
Substituting $\tilde K$ into $Q+\Acl^\top X\Acl+\tilde K^\top\tilde R\tilde K$ and
cancelling yields
\begin{equation}\label{eq:garelyap}
X=Q+\tilde K^\top\tilde R\tilde K+\Acl^\top X\Acl,
\end{equation}
\emph{identical algebra} to the LQG identity below \eqref{eq:cos}; completion of
squares never uses the sign of $\tilde R$, which is why indefiniteness is harmless.

\paragraph{(c) Local uniqueness and $C^1$ (the heart).}
The GARE residual $\mathcal R(X,\Th)$ (right minus left of \eqref{eq:gare}) is
$C^1$ wherever the pivot is nonsingular, hence near a stabilizing root. Because
$\tilde K$ is the \emph{saddle point} of the inner game (nonsingular pivot), the
game envelope annihilates the gain's first-order variation and the partial
derivative in $X$ collapses to the Stein operator
\begin{equation}\label{eq:gareift}
\partial_X\mathcal R(X,\Th)[\delta X]=\Acl^\top\delta X\,\Acl-\delta X,
\end{equation}
invertible precisely because $\Acl$ is Schur. The \emph{implicit function theorem}
then gives a unique $C^1$ branch $\Th\mapsto X(\Th)$ through the stabilizing root,
and with it $C^1$ maps $\tilde K(\Th),\Acl(\Th)$ and the game gramian $S(\Th)$
solving $S=\Acl S\Acl^\top+W_\Th$. This is Proposition~\ref{prop:hinf}; note it
needs no dual and no strict complementarity --- the regularity is read straight off
the Riccati map, unlike the SDP route of Section~\ref{sec:lqg}.

\paragraph{(d) $C^1$ value and gradient.}
$\Vwc=\tr(XW_\Th)$ is $C^1$ by (c) and (S1). Differentiating \eqref{eq:garelyap},
saddle stationarity of $\tilde K$ kills all $\mathrm d\tilde K$ terms, so
$\mathrm dX$ solves the Stein equation
$\mathrm dX-\Acl^\top\mathrm dX\,\Acl=\Xi$ with source
$\Xi=\mathrm dQ+\tilde K^\top\mathrm d\tilde R\tilde K
+\Acl^\top X(\mathrm dA_\Th+\mathrm d\tilde B_\Th\tilde K)+(\cdot)^\top$.
Contracting with the gramian and using
$\tr(W_\Th\,\mathrm dX)=\tr\!\big(S(\mathrm dX-\Acl^\top\mathrm dX\,\Acl)\big)
=\tr(S\Xi)$ gives, term by term,
\begin{equation}\label{eq:garegrad}
\frac{\partial\Vwc}{\partial A_\Th}=2X\Acl S,\qquad
\frac{\partial\Vwc}{\partial\tilde B_\Th}=2X\Acl S\tilde K^\top,\qquad
\frac{\partial\Vwc}{\partial Q}=S,\qquad
\frac{\partial\Vwc}{\partial\tilde R}=\tilde KS\tilde K^\top,
\end{equation}
plus the explicit $\partial\Vwc/\partial W_\Th=X$ from the objective. The channel
$E_\Th=W_\Th^{\half}$ enters through the $E$-block of $\partial\Vwc/\partial\tilde
B_\Th$, so a design coordinate moving both the objective weight $W_\Th$ and the
channel $E_\Th$ sums the two contributions (the data-object chain rule of
Assumption~1(A4)).

\paragraph{(e) Filter GARE by duality; (f) the $H_2$ limit.}
Under the transpose substitution $A_\Th\mapsto A_\Th^\top$,
$\tilde B_\Th\mapsto[\,C_\Th^\top\ L_e^\top\,]$,
$\tilde R\mapsto\diag(V,-\gamma^2I)$, the worst-case error covariance $\Sigma_e$
solves the control GARE on the dual data, so (a)--(d) transfer and
$\Vest=\tr(M\Sigma_e)\in C^1$ with the dual gradient. Finally, as
$\gamma\to\infty$ the $-\gamma^2I$ pivot block deactivates, \eqref{eq:gare} reduces
to the control DARE, $X\to P$, $S\to\Sigma^\star$, and \eqref{eq:garegrad}
collapses to the LQG gradients of Section~\ref{sec:lqg}: the $H_2$/LQG result is the
$\gamma\to\infty$ boundary of this one.

%=======================================================================
\section{What each proof rests on}\label{sec:summary}
%=======================================================================

\begin{itemize}
\item \textbf{Slater (LQG only)} $\Leftarrow$ a stabilizing feedback plus an
$\varepsilon I$ inflation, i.e.\ the controllability/observability data, with
$W,V\succ0$ supplying the strict margin; conic strong duality then gives (A1)
\cite{BonnansShapiro,BoydVandenberghe}.
\item \textbf{Primal uniqueness} $\Leftarrow$ the (game) Riccati equation: the
covariance/certificate SDP is an exact convex lift whose optimizer is the
(G)ARE solution \cite{LancasterRodman,AndersonMoore,BasarBernhard}, pinned by
completion of squares.
\item \textbf{Differentiability and gradient correctness} $\Leftarrow$ the
\emph{implicit function theorem} through the Stein operator ($\Acl$ Schur), which
alone makes the value $C^1$; and, in the SDP case, the
\emph{Danskin/Bonnans--Shapiro} formula, which makes the envelope gradient
\emph{dual-invariant} --- so the returned-dual gradient is exact even without
strict complementarity \cite{BonnansShapiro}.
\item \textbf{Dual uniqueness (A3)} $\Leftarrow$ strict complementarity /
nondegeneracy in the \emph{Alizadeh--Haeberly--Overton} sense
\cite{AHO,Shapiro97}; needed only for the ``unique dual'' upgrade, not for a
correct gradient.
\item \textbf{$H_\infty$ vs.\ LQG.} The $H_\infty$ proof is the Riccati half only:
IFT $+$ Stein operator $+$ gradient, with an indefinite weight and no SDP dual. It
therefore never needs Slater, dual invariance, or strict complementarity --- the
fixed-$\gamma$ GARE is well-posed by construction, and the $H_2$ case is its
$\gamma\to\infty$ limit.
\end{itemize}

%=======================================================================
\begin{thebibliography}{9}
%=======================================================================

\bibitem{AHO}
F.~Alizadeh, J.-P.~A.~Haeberly, and M.~L.~Overton,
\emph{Complementarity and nondegeneracy in semidefinite programming},
Math.\ Program.\ \textbf{77} (1997), 111--128.

\bibitem{Shapiro97}
A.~Shapiro,
\emph{First and second order analysis of nonlinear semidefinite programs},
Math.\ Program.\ \textbf{77} (1997), 301--320.

\bibitem{BonnansShapiro}
J.~F.~Bonnans and A.~Shapiro,
\emph{Perturbation Analysis of Optimization Problems},
Springer, New York, 2000.

\bibitem{BoydVandenberghe}
S.~Boyd and L.~Vandenberghe,
\emph{Convex Optimization},
Cambridge University Press, 2004.

\bibitem{BoydLMI}
S.~Boyd, L.~El~Ghaoui, E.~Feron, and V.~Balakrishnan,
\emph{Linear Matrix Inequalities in System and Control Theory},
SIAM, Philadelphia, 1994.

\bibitem{AndersonMoore}
B.~D.~O.~Anderson and J.~B.~Moore,
\emph{Optimal Control: Linear Quadratic Methods},
Prentice Hall, Englewood Cliffs, NJ, 1990.

\bibitem{LancasterRodman}
P.~Lancaster and L.~Rodman,
\emph{Algebraic Riccati Equations},
Oxford University Press, 1995.

\bibitem{BasarBernhard}
T.~Ba\c{s}ar and P.~Bernhard,
\emph{$H_\infty$-Optimal Control and Related Minimax Design Problems:
A Dynamic Game Approach}, 2nd ed.,
Birkh\"auser, Boston, 1995.

\end{thebibliography}

\end{document}
